\documentclass[10pt,a4paper]{amsart}
\usepackage[margin=19mm]{geometry}
\usepackage{amsmath,amssymb,mathtools}
\usepackage[scr=rsfs,cal=euler]{mathalfa}
\usepackage{xfrac}
\usepackage[unicode,hidelinks,bookmarks=false]{hyperref}
\newcommand{\ie}{i.e.\ }
\newcommand{\cf}{cf.\ }
\newcommand{\resp}{respectively\ }
\newcommand{\Subsection}{\subsection}
\providecommand{\nobreakdash}{\nobreak\hbox{-}}
\setlength{\emergencystretch}{2em}
\multlinegap0pt
\usepackage{amssymb}
\usepackage{bm}
\usepackage[shortlabels]{enumitem}
\newtheorem{prop}{Proposition}
\newtheorem{lemma}{Lemma}
\def\D{\textup{d}} 
\renewcommand{\Im}{\operatorname{Im}}
\def\R{\mathbb{R}}
\def\ad{\operatorname{ad}}
\def\g{\mathfrak{g}}
\def\h{\mathfrak{h}}
\def\pint{\langle \cdotp,\cdotp \rangle }
\def\q{\mathfrak{q}}
\def\rad{\operatorname{rad}}
\newcommand{\restr}{\raisebox{-.2em}{$\!\big\arrowvert$}}
\def\s{\mathfrak s}
\newcommand{\then}{\Longrightarrow}
\title{$\boldsymbol \eta$-Einstein Sasakian Lie algebras}
\author{Adri{\'a}n Andrada, Simon G. Chiossi, Alberth J. N{\'u}\~nez Sullca}
\date{Source: arXiv:2504.16033v2 (CC BY 4.0)}
\begin{document}
\maketitle

\noindent\emph{Source and licence.} $\boldsymbol \eta$-Einstein Sasakian Lie algebras. Version arXiv:2504.16033v2 (CC BY 4.0), distributed by arXiv under the Creative Commons Attribution 4.0 International licence, http://creativecommons.org/licenses/by/4.0/, as recorded in the arXiv licence metadata for this exact version and shown on the abstract page https://arxiv.org/abs/2504.16033v2. Full licence notice: \texttt{licenses/arxiv-2504.16033-CC-BY-4.0.txt}.\par\smallskip

\noindent\textit{Editorial scope.} Counted unit: the article's lemma radical (source0.tex lines 668--671). The article's complete original proof of it is delivered with it (source0.tex lines 673--743, one \texttt{proof} environment of 71 lines). No mathematics is written, repaired or re-derived: the statement and the proof are the article's own lines, in the article's own notation, and no retained line is substituted or re-flowed.  Retained uncounted as the source-local context the proof needs: lines 321--321; lines 487--489; lines 507--510; lines 512--516; lines 519--521; lines 574--589.  Nothing was omitted from any retained span, so the package carries no omission record.  No retained line contains a citation, so the wrapper carries no bibliography and no bibliography entry.  No retained line contains a figure environment, an image-inclusion command or drawing code, so no image is delivered and the package declares no asset dependency.  Editorial formatting only, in addition: an AMS \texttt{amsart} preamble that restates the article's own declarations and macros used by the retained lines, namely the environments \texttt{prop}, \texttt{lemma} on separate counters, together with the standard packages those lines need.  The retained environments print in this document as Proposition 1, Lemma 1 in order of appearance, whereas the article prints the same items with the numbering of its own full document.  The complete native source of the article as distributed in the arXiv e-print for this exact version is bundled unchanged as \texttt{sources/176917/source0.tex}.
\section{Einstein Sasakian algebras with centre}\label{sec:ES with centre} 

\begin{equation}\label{eq: h1-bracket}
 [x,y]=-2\omega(x,y)\xi + \theta(x,y), \quad [x,\xi]=0, \qquad x,y \in \h_1.
\end{equation} 

\begin{equation}\label{eq:ad_xi w}
\ad_{\xi}\restr_{\q}=
\begin{pmatrix} 0 & -k \\ k & 0 \end{pmatrix}
\end{equation}

\begin{equation}\label{eq:ad_x w}
\ad_{x}\restr_{\q}=
\begin{pmatrix} \mu(x) & -\gamma(x) \\ \gamma(x) & \mu(x)
\end{pmatrix},\qquad x\in \h_1.
\end{equation}

\begin{equation} \label{[u,v]} 
    [u,v]=2\xi + H_{0}\quad \text{ for some }\ H_{0} \in \h_1.
\end{equation}

\begin{prop}\label{h1 crea sasaki sin centro}
    Let  $(\h_1, \theta(\cdot,\cdot), \pint)$ be a metric Lie algebra of dimension $2m$  equipped with a K{\"a}hler-exact structure $\omega=-\frac{1}{2k}\D^{\h_1}\gamma$, for some $k\neq 0$, $\gamma \in \h_{1}^*$.
 Assume there exist an element $H_0 \in \theta(\h_1,\h_1)$ obeying the compatibility conditions:
\begin{align}
    & JH_0 \in \theta(\h_1,\h_1)^{\perp}  \label{eq: JH0}\\ 
    & 2\mu(x)H_0=\theta(x,H_0) \quad \text{for every } x \in \h_1, \label{eq: importante}
\end{align}
where $\mu\in\h_1^*$ is the closed 1-form given by $\mu(x)=-\frac12\langle x,JH_0\rangle$, $x\in \h_1$. \par 
Call $\s= \R\xi \oplus \h_1$ the (Sasakian) central extension of $\h_1$ by the cocycle $\omega$, as described in \S \ref{sec:ES with centre}.  
Choose a $2$-plane ${\q}$ with orthonormal basis $(u,v)$ and let 
\[ \g=\s \oplus {\q}\]
with bracket prescribed by \eqref{eq: h1-bracket}, \eqref{eq:ad_xi w}, \eqref{eq:ad_x w} and \eqref{[u,v]}. This makes  $\g$ a Lie algebra of dimension $2m+3$ with $\s$ as subalgebra. \par
Moreover, if we extend the structure $(\varphi, \eta, \xi, \pint)$ on $\s$ to $\g$ by 
$$\varphi (u)=v,\qquad \eta({\q})=0,$$  
and the rest in the natural way, $(\g, \varphi, \eta, \xi , \pint)$ becomes a Sasakian Lie algebra with no centre.
\end{prop}

\begin{lemma}\label{lemma: radical}
Let $\g$ be a Sasakian Lie algebra with trivial centre and $\dim \Im \ad_{\xi}=2$. If $H_0=0$ then 
$\rad(\g)= \left\{ -\dfrac{\gamma(h)}{k} \xi+h \mid  h\in \rad(\h_1) \right\}.$
\end{lemma}

\begin{proof}
Let us recall that 
\begin{equation*}\label{radical formula}
       \rad(\g)=[\g,\g]^{\perp _{\beta}}=\{z \in \g \mid \beta(z, y)=0 \, \text{ for all}\; y \in [\g,\g]\},
\end{equation*}
where $\beta$ is the Cartan-Killing form of $\g$.  
 Having $H_0=0$ implies $\mu=0$ and $[u,v]=2\xi$, so by Proposition \ref{h1 crea sasaki sin centro} 
\begin{equation}\label{conmutador de no soluble}
    [\g,\g]=\R\xi \oplus [\h_1, \h_1] \oplus {\q}. 
\end{equation}
Let us, for clarity, write the matrices used to compute $\beta$ in the orthonormal frame $(\xi, e_1,\ldots, e_{2n},u,v)$ of $\g$, where $(e_i)$ is an orthonormal basis of $\h_1$. 
The matrix of $\ad_{\xi}$ is all zero except for the diagonal ${\q}$-block where 
\eqref{eq:ad_xi w} holds.
For any $h\in\h_1$, 
\[
\ad_{h}=
\left(\begin{array}{c|ccc|cc}
          0  &\cdots & 2\langle Jh,e_i\rangle & \cdots &0&0 \\
          \hline 
          0&&& &0&0\\
          \vdots&&\ad^{\h_1}_h&&\vdots&\vdots\\
          0&&&&0&0\\
          \hline
       0 & \cdots & 0 & \cdots &0 & -\gamma(h) \\
        0 & \cdots & 0 & \cdots &\gamma(h) & 0
        \end{array}\right).
\]
Moreover,
\[
\ad_{u}=
\left(\begin{array}{c|ccc|cc}
          0  &\cdots & 0 & \cdots&0&2 \\
          \hline 
          0&&& &0&0\\
          \vdots&&0&&\vdots&\vdots\\
          0&&&&0&0\\
          \hline
       0 & \cdots & 0 & \cdots&0 & 0 \\
        -k &  \cdots & -\gamma(e_i) &  \cdots
        &0 & 0
        \end{array}\right), \quad 
\ad_{v}=
\left(\begin{array}{c|ccc|cc}
          0  &\cdots & 0 & \cdots&-2&0 \\
          \hline 
          0&&& &0&0\\
          \vdots&&0&&\vdots&\vdots\\
          0&&&&0&0\\
          \hline
       k &  \cdots & \gamma(e_i) &  \cdots
       &0 & 0 \\
        0 & \cdots & 0 & \cdots&0 & 0
        \end{array}\right).
\]
This gives the values of $\beta$  (for $h,\tilde{h}\in\h_1$):
$$\begin{array}{ll}
\beta(\xi,\xi)=-2k^2, & \; \beta(u,u)=-4k, \qquad \beta(v, v)=-4k\\
\beta(\xi, h)=-2k\gamma(h), & \;
\beta(h,\widetilde{h})=\beta^{\h_1}(h,\widetilde{h})-2\gamma(h)\gamma(\widetilde{h}).
\end{array}$$

Pick an element $z=a\xi+ bu +cv+h \in \rad(\g)$, where $a,b,c\in\R$, $h\in\h_1$. Using decomposition \eqref{conmutador de no soluble} we obtain
\begin{equation*}\begin{split}
    0=\beta(z,\xi)&=-2ak^{2}-2k\gamma(h) \then \gamma(h)=-ak \\
    0=\beta(z,u)&=-4b k \then b=0\\
    0=\beta(z, v)&=-4ck \then c=0.
\end{split}\end{equation*}
Consequently for $h' \in [\h_1, \h_1]$
$$0=\beta(z, h')=-2ak\gamma(h')+\beta^{\h_1}(h, h') -2\gamma(h)\gamma(h') = \beta^{\h_1}(h, h'),$$
which says $h \in \rad(\h_1)$.
\end{proof}

\end{document}
